% SPDX-License-Identifier: Apache-2.0
\section{A Unit That Answers the Reset Writes Its Memory}
\label{sec:x-rstmemown}

\S\ref{sec:x-rstmem} drops a memory write made while the reset is
held, and \S\ref{sec:x-ownrst} leaves a unit that declares
\code{rst} to say in its body what the reset does to its registers.
The two meet in a unit that declares \code{rst} and writes a memory:
its memory is its body's too. The netlist gives such a unit no branch
out of reset, so its writes are not gated, and the run keeps them as
well (issue 966).

The clearest reason to want it is \code{Wipe}: a store that clears
itself while the reset is held, a word a cycle, so that it comes out
of reset empty whatever it held before. The run fills the store, holds
the reset for as many cycles as the store has words, then writes one
word and reads all four back. Every word is the one the wipe wrote,
except the one written after it. A memory is not traced, so the
simulations of the netlist see the store through \code{q} and are
checked against the run there.

\begin{figure}[htbp]
\centering
\input{rstmemown_timing.tex}
\caption{The store filled, four cycles of \code{rst} with \code{at}
walking the words, one write at word 2, and the store read back:
\code{0xee} everywhere but word 2.}
\label{fig:x-rstmemown}
\end{figure}

\lstinputlisting[caption={\code{ex\_rstmemown.rs}. Run with \code{bazel run //lib/examples:ex\_rstmemown}.},label={lst:x-rstmemown}]{lib/examples/ex_rstmemown.rs}

\lstinputlisting[style=ascii,caption={What \code{ex\_rstmemown} printed when the build ran it: the store read back, then the netlist, whose memory write is not gated on the reset.},label={lst:x-rstmemown-out}]{out_rstmemown.txt}
